% The storefront query plan at tag ch11 and at tag ch12. Referenced from
% chapter 12. Bare tikzpicture; the figure environment, caption and label live
% at the call site.
%
% The point of the picture is the middle node on the right. Three of the four
% new fetches happen at once because they depend on the same key and on nothing
% else, and the fourth cannot, because its keys do not exist until the fetch
% above it has answered. Depth in the graph is sequencing in the plan.
\begin{tikzpicture}[
    font=\small,
    head/.style={font=\small\bfseries, align=center},
    fetch/.style={draw, rounded corners=2pt, minimum width=20mm,
                  minimum height=5.5mm, align=center, font=\scriptsize},
    par/.style={draw, dashed, rounded corners=2pt, gray},
    flow/.style={-{Stealth[length=2mm]}},
    note/.style={font=\scriptsize, gray, align=center, inner sep=1pt}
  ]

  % -- left: two fetches ------------------------------------------------------
  \node[head] at (1.5,0.8) {two, at \code{ch11}};

  \node[fetch] (a1) at (1.5,0.0)  {\code{catalog}};
  \node[fetch] (a2) at (1.5,-1.5) {\code{mosaic}};

  \draw[flow] (a1) -- (a2);
  \node[note, anchor=west] at (2.7,-0.75) {one key};

  % -- right: five fetches in three steps -------------------------------------
  \node[head] at (8.4,0.8) {five, at \code{ch12}};

  \node[fetch] (b1) at (8.4,0.0)   {\code{catalog}};

  \draw[par] (5.6,-1.9) rectangle (11.2,-1.1);
  \node[fetch] (bp) at (6.7,-1.5)  {\code{pricing}};
  \node[fetch] (bi) at (8.4,-1.5)  {\code{inventory}};
  \node[fetch] (br) at (10.1,-1.5) {\code{reviews}};

  \node[fetch] (ba) at (10.1,-3.1) {\code{accounts}};

  \draw[flow] (b1) -- (bp);
  \draw[flow] (b1) -- (bi);
  \draw[flow] (b1) -- (br);
  \draw[flow] (br) -- (ba);

  \node[note, anchor=east, align=right] at (5.5,-1.5)
    {at once: same\\key, nothing else};
  \node[note, anchor=east, align=right] at (8.9,-3.1)
    {after: these keys do not\\exist until reviews answers};

\end{tikzpicture}
